% Source: Anteproyecto (3).pdf. Original wording retained except population/sample.
\chapter{Diseño metodológico}

El presente diseño metodológico establece la ruta estructurada que orientará el desarrollo del proyecto desde la identificación de los datos hasta el despliegue del prototipo funcional. El proyecto se enmarca en una investigación aplicada con enfoque cuantitativo y alcance experimental, cuyo propósito es el diseño e implementación de un modelo de inteligencia artificial basado en aprendizaje autosupervisado para el apoyo diagnóstico de lesiones cutáneas en contextos clínicos con escasez de datos etiquetados en el departamento de Santander.

El desempeño del modelo será cuantificado mediante métricas estándar de clasificación AUC-ROC, F1-score y exactitud que permiten comparaciones objetivas entre configuraciones (Hernández-Sampieri et al., 2018). Para garantizar la validez interna de las comparaciones, el modelo SSL será entrenado y sometido a ajuste de hiperparámetros manteniendo constantes el conjunto de datos, la arquitectura backbone y el protocolo de evaluación, contrastando los resultados obtenidos con modelos supervisados de referencia reportados en la literatura bajo las mismas condiciones experimentales.

\section{Diseño de la investigación}

El desarrollo del proyecto estará basado en el cumplimiento de los objetivos formulados siguiendo el modelo CRISP-ML(Q), estructurado originalmente en seis fases: comprensión del negocio y de los datos, ingeniería de datos, ingeniería del modelo, evaluación del modelo, despliegue y monitoreo continuo, Estas fases se agrupan en cuatro etapas iterativas representadas en la Figura~\ref{fig:original-10}: (i) Comprensión del problema clínico y de los datos dermatológicos, (ii) Ingeniería de datos y modelado del algoritmo SSL, (iii) Evaluación del modelo con aseguramiento de calidad y (iv) Implementación y despliegue del prototipo de apoyo al diagnóstico clínico. La sexta fase del modelo original, correspondiente al monitoreo y mantenimiento del modelo en producción, excede el nivel de madurez tecnológica TRL 4 proyectado para este trabajo de grado; por esta razón, se documenta como línea de continuidad para investigaciones posteriores.

\begin{figure}[htbp]
    \centering
    \includegraphics[width=0.9\textwidth]{figura10-diseno-crisp-mlq.png}
    \caption{Diseño metodológico basado en CRISP-ML(Q)}
    \label{fig:original-10}
\par\smallskip
\raggedright\footnotesize
\textit{Nota.} Creación propia.
\end{figure}

\section{Población y muestra}

La disponibilidad de conjuntos de datos de imágenes dermatológicas con fines investigativos es reducida en el contexto clínico del departamento de Santander. Ante esta restricción, el estudio adopta repositorios públicos consolidados en la literatura de clasificación dermatológica. La población de análisis corresponde a las imágenes dermatológicas disponibles en los conjuntos seleccionados; su uso no constituye una muestra representativa de pacientes del departamento ni una validación clínica regional.

HAM10000 es el conjunto fuente para el entrenamiento del clasificador y para la selección, calibración y evaluación interna, con sus siete categorías diagnósticas \parencite{tschandl2018ham10000}. Las particiones deben mantener cada lesión identificada por \texttt{lesion\_id} dentro de un único subconjunto para prevenir la fuga entre entrenamiento, validación y prueba. Los tamaños efectivos de la muestra y de las particiones se documentarán mediante los manifiestos de cada experimento, sin asumir que todas las imágenes disponibles se utilizan en cada fase.

PAD-UFES-20 se reserva para la evaluación externa del modelo \parencite{pacheco2020padufes}. Sus imágenes clínicas capturadas con teléfonos móviles difieren de las imágenes dermatoscópicas de HAM10000, de modo que esta evaluación estudia la transferencia entre modalidades y dominios. PAD-UFES-20 no se utiliza para el entrenamiento, la selección de modelos o hiperparámetros, la calibración ni el ajuste de umbrales. Se documentarán el mapeo de etiquetas, las clases compatibles y las limitaciones de la comparación.

La selección de imágenes estará guiada por criterios de calidad y pertinencia diagnóstica, aplicados de forma diferenciada según la fuente. La procedencia, las etiquetas disponibles, las licencias, los criterios de inclusión y exclusión y el preprocesamiento se registrarán junto con los identificadores y hashes de las particiones. La composición del conjunto de imágenes no etiquetadas para preentrenamiento se reportará conforme a los manifiestos y exclusiones de cada corrida, evitando incorporar imágenes reservadas para selección o prueba.

\section{Fases de la investigación}

\subsection{Fase 1 -- Comprensión del problema clínico y de los datos dermatológicos}

En esta fase se desarrolla la comprensión integral del problema diagnóstico de lesiones cutáneas en el contexto del departamento de Santander y se realiza la caracterización de los conjuntos de datos públicos de imágenes dermatológicas que sustentarán el desarrollo del modelo. Se evalúan aspectos relacionados con la incidencia y prevalencia epidemiológica de las lesiones cutáneas reportadas a nivel regional, las propiedades estadísticas de los datasets candidatos --distribución de clases, calidad de imagen, tamaño y condiciones de licencia-- y la matriz de riesgos por fase que orientará las etapas posteriores del proyecto, en concordancia con los componentes de aseguramiento de calidad (Q) propuestos por Studer et al. (2021). Entregable: Informe de comprensión empresarial y de datos con la caracterización formal de los conjuntos de datos seleccionados, los criterios de inclusión y exclusión, las métricas objetivo del modelo y la matriz de riesgos por fase.

Actividad 1.1: Revisar literatura científica y reportes epidemiológicos relacionados con lesiones cutáneas de mayor relevancia clínica en el departamento de Santander.

Actividad 1.2: Definir los objetivos de aprendizaje automático del proyecto y traducirlos a métricas medibles, en concordancia con el perfil epidemiológico regional y los estándares reportados en la literatura.

Actividad 1.3: Identificar y caracterizar conjuntos de datos públicos de imágenes dermatológicas, documentando tamaño, calidad de imagen, distribución de clases, fototipos representados y condiciones de licencia.

Actividad 1.4: Seleccionar las imágenes y clases diagnósticas pertinentes para el contexto del proyecto, aplicando criterios de inclusión y exclusión previamente definidos.

Actividad 1.5: Consolidar la matriz de riesgos por fase y los criterios de aseguramiento de calidad (Q) que orientarán el desarrollo de las fases posteriores.

\subsection{Fase 2 -- Ingeniería de datos y modelado del algoritmo SSL}

En esta fase se desarrolla el pipeline reproducible de preparación de los datos seleccionados en la fase anterior y se diseña el algoritmo de clasificación basado en aprendizaje autosupervisado. El proceso incluye la implementación de las transformaciones de preprocesamiento, la partición estratificada de los conjuntos de datos, la selección de la arquitectura backbone, el entrenamiento autosupervisado con imágenes no etiquetadas y el ajuste fino supervisado, manteniendo un registro completo de metadata de cada experimento para garantizar la reproducibilidad de los resultados. Entregable: Pipeline de datos versionado y algoritmo computacional basado en aprendizaje autosupervisado para la clasificación de lesiones cutáneas, con documentación de arquitectura, hiperparámetros y trazabilidad experimental.

Actividad 2.1: Implementar el pipeline de preprocesamiento de imágenes, incluyendo redimensionamiento, normalización, aumento de datos y partición estratificada en subconjuntos de entrenamiento, validación y prueba.

Actividad 2.2: Revisar modelos y estrategias de aprendizaje autosupervisado reportadas en la literatura para clasificación de imágenes médicas y dermatológicas.

Actividad 2.3: Seleccionar la arquitectura base (backbone) y definir la estrategia de entrenamiento autosupervisado del modelo propuesto.

Actividad 2.4: Entrenar el modelo mediante aprendizaje autosupervisado utilizando imágenes dermatológicas sin etiquetas y realizar el ajuste fino supervisado sobre el conjunto de datos etiquetados.

Actividad 2.5: Documentar la arquitectura final del modelo, los hiperparámetros utilizados, las decisiones de diseño adoptadas y la metadata completa de los experimentos realizados.

\subsection{Fase 3 -- Evaluación del modelo con aseguramiento de calidad}

En esta fase se evalúa el desempeño del algoritmo propuesto mediante métricas estándar de clasificación dermatológica reportadas en la literatura científica y se realiza una comparación experimental con modelos supervisados de referencia entrenados bajo las mismas condiciones experimentales. Adicionalmente, se incorporan los componentes de aseguramiento de calidad (Q) propios del marco CRISP-ML(Q), incluyendo el análisis de robustez del modelo, la auditoría de fairness por fototipo de piel y la generación de mapas de explicabilidad. Entregable: Informe comparativo con métricas de desempeño, matrices de confusión, análisis estadístico, auditoría de fairness y mapas de explicabilidad del modelo propuesto frente a los modelos supervisados de referencia.

Actividad 3.1: Identificar y seleccionar modelos supervisados de referencia reportados en la literatura científica para comparación experimental.

Actividad 3.2: Implementar los modelos de referencia utilizando el mismo conjunto de datos, backbone y protocolo experimental del modelo propuesto.

Actividad 3.3: Calcular métricas de evaluación como AUC-ROC, F1-score, exactitud, precisión y sensibilidad, desagregando los resultados por clase diagnóstica.

Actividad 3.4: Realizar análisis de robustez del modelo, auditoría de fairness por fototipo de piel y generación de mapas de explicabilidad como componentes de aseguramiento de calidad.

Actividad 3.5: Elaborar el análisis estadístico y la documentación comparativa de los resultados obtenidos durante la evaluación experimental.

\subsection{Fase 4 -- Implementación y despliegue del prototipo de apoyo al diagnóstico clínico}

En esta fase se desarrolla un prototipo funcional que integra el algoritmo de clasificación validado para apoyar el análisis de imágenes dermatológicas en un entorno clínico experimental. El aplicativo permite cargar imágenes dermatológicas, visualizar las predicciones generadas por el modelo, declarar de forma explícita las limitaciones conocidas del sistema y aplicar un plan de contingencia frente a casos de baja confianza, en concordancia con los principios de despliegue responsable propuestos por el marco CRISP-ML(Q). Entregable: Prototipo funcional documentado con manual de uso, matriz de limitaciones y plan de contingencia.

Actividad 4.1: Definir los requerimientos funcionales y no funcionales del aplicativo de apoyo diagnóstico.

Actividad 4.2: Seleccionar las herramientas y recursos tecnológicos requeridos para el desarrollo del prototipo, priorizando tecnologías de acceso abierto.

Actividad 4.3: Diseñar e implementar la interfaz del aplicativo para la carga de imágenes y la visualización de la predicción del modelo junto con sus métricas de desempeño y limitaciones.

Actividad 4.4: Integrar el modelo entrenado dentro del flujo funcional del aplicativo y configurar el plan de contingencia frente a casos de baja confianza.

Actividad 4.5: Realizar pruebas funcionales y de validación del prototipo bajo diferentes condiciones experimentales, verificando la consistencia de las predicciones generadas por el sistema.

\section{Análisis de factibilidad}

El proyecto presenta viabilidad técnica al disponer de las arquitecturas SSL candidatas SimCLR, BYOL, DINO y MAE como implementaciones de referencia en repositorios públicos bajo licencias de código abierto, complementadas con frameworks consolidados en la comunidad de aprendizaje profundo (PyTorch, scikit-learn) y con los conjuntos de datos HAM10000 e ISIC Archive disponibles bajo licencias académicas para uso investigativo, configuración que proyecta un nivel de madurez tecnológica TRL 4 al cierre del trabajo. A nivel computacional, las fases de análisis exploratorio y caracterización de los conjuntos de datos se ejecutan en Google Colab Pro y Kaggle Notebooks, aprovechando sus GPUs NVIDIA T4 y A100 de acceso gratuito; el entrenamiento del modelo se desarrollará sobre los recursos de cómputo propios del equipo, con la articulación potencial de la infraestructura institucional de la Universidad Autónoma de Bucaramanga conforme a su disponibilidad durante el cronograma; y el producto final un prototipo dirigido a médicos generales y dermatólogos. En la dimensión operativa, la co-dirección de Karen Yaneth Sánchez Quiroga, investigadora afiliada al King Abdullah University of Science and Technology (KAUST), aporta experiencia directa en sistemas de diagnóstico asistido por computador (CAD) sobre el dataset HAM10000 (Calderón, Sánchez \& Argüello, 2021) y la autoría del dataset CO2Wounds-V2 recolectado en Contratación, Santander (Sánchez et al., 2024).

\section{Análisis de riesgos y planes de contingencia}

El presente apartado consolida los principales riesgos identificados durante la planificación del proyecto, así como sus respectivos planes de contingencia, en concordancia con la naturaleza iterativa de la metodología CRISP-DM (Wirth \& Hipp, 2000). En la Tabla~\ref{tab:riesgos-proyecto} se presentan siete riesgos clasificados por nivel de criticidad (Alto o Medio) y se describe, para cada uno, una estrategia de mitigación y el punto de retorno previsto dentro del ciclo metodológico cuando el riesgo se materialice. Los tres primeros riesgos (R1, R2 y R3) corresponden a las amenazas técnicas inherentes al desarrollo de modelos de aprendizaje automático con datos clínicos limitados, mientras que los riesgos R4 a R7 abordan dimensiones transversales del proyecto, tales como el sesgo demográfico de los conjuntos de datos públicos, la cobertura epidemiológica regional, la trazabilidad experimental y la consistencia entre el modelo entrenado y su despliegue en el prototipo de aplicación. Esta organización permite anticipar con criterio metodológico los escenarios que podrían comprometer el cumplimiento de los objetivos y articular respuestas verificables que mantengan la coherencia entre las fases del proyecto.

\begingroup
\small
\singlespacing
\setlength{\tabcolsep}{4pt}
\begin{longtable}{@{}>{\raggedright\arraybackslash}p{0.34\textwidth}>{\raggedright\arraybackslash}p{0.09\textwidth}>{\raggedright\arraybackslash}p{0.51\textwidth}@{}}
\caption{Riesgos del proyecto y planes de contingencia}\label{tab:riesgos-proyecto}\\
\toprule
\textbf{Riesgo} & \textbf{Nivel} & \textbf{Plan de contingencia} \\
\midrule
\endfirsthead
\toprule
\textbf{Riesgo} & \textbf{Nivel} & \textbf{Plan de contingencia} \\
\midrule
\endhead
\bottomrule
\endfoot
R1. Desempeño insuficiente del modelo SSL. & Alto & Explorar arquitecturas SSL alternativas y modificar la estrategia de aumentación. Retorno: Fase 3 - Fase 2. \\
\addlinespace
R2. Desbalance severo de clases en el dataset. & Alto & Aplicar sobremuestreo (SMOTE) o pesos por clase en la función de pérdida. Retorno: Fase 2. \\
\addlinespace
R3. Limitaciones de cómputo (memoria/GPU). & Medio & Adoptar arquitecturas más ligeras como backbone, priorizando la viabilidad. Retorno: Fase 2. \\
\addlinespace
R4. Sesgo de fototipos IV-V en conjuntos de datos públicos. & Medio & Documentar como limitación; reportar desempeño desagregado y proponer trabajo futuro con datos regionales. \\
\addlinespace
R5. Cobertura de clases vs. perfil epidemiológico regional. & Medio & Combinar varios datasets compatibles; ajustar criterios de inclusión. Retorno: Fase 1. \\
\addlinespace
R6. Pérdida de trazabilidad de experimentos. & Medio & Registrar cada experimento en plataforma de seguimiento; versionar el código con Git. Control transversal. \\
\addlinespace
R7. Discrepancia laboratorio vs. prototipo. & Medio & Pruebas funcionales con imágenes externas; alinear preprocesamiento. Retorno: Fase 4 - Fase 3. \\
\addlinespace
\end{longtable}
\endgroup
 Nota. Elaboración propia con base en el diseño metodológico del proyecto (apartado 7.6). La columna "Retorno CRISP-DM" embebida en el plan de contingencia indica la fase a la que se regresa dentro del ciclo iterativo (Wirth \& Hipp, 2000) cuando el riesgo se materializa. Convención de niveles: Alto, Medio, Bajo.
